\begin{abstract}
  Model merging has emerged as a key paradigm for combining task-specific expert checkpoints into a single unified model without retraining. In this work, we present \textbf{FluidMerge} (Fluid-Dynamic Parameter Coalescence) as an intuitive and interesting physical analog that views the parameter trajectory during test-time adaptation as a continuous-time advection-diffusion fluid flow. To overcome representation misalignment and the coordinate-dependency of spatial grid Laplacians, we design, implement, and evaluate two core techniques: \textbf{Expert-Weighted Initial Boundary Conditions} (which are mathematically equivalent to initializing optimization from the standard Task Arithmetic weight-space average of Ilharco et al., 2023) and \textbf{Fisher-Information-based Viscosity} (which, under discretized Euler integration, is mathematically isomorphic to a continuous-time formulation of Elastic Weight Consolidation of Kirkpatrick et al., 2017). We conduct a rigorous empirical evaluation of FluidMerge and state-of-the-art post-hoc adaptation baselines across eight image classification datasets using a Vision Transformer (ViT-B-32). Our results demonstrate that starting from the raw base weights leads to a formidable \emph{domain shift barrier} (random-guess performance of $\sim 5\%$), whereas initializing from the Task Arithmetic boundary condition places the parameters within high-performing multi-task basins. We show that Fisher-Information viscosity (acting as a point-wise, function-sensitive restorative force) successfully prevents representation collapse and stabilizes the continuous-time parameter trajectory, outperforming standard $L_2$ weight anchoring and baseline methods. This work establishes a clear, transparent, and physically motivated framework for continuous-time parameter adaptation.
\end{abstract}
